\section{多模态语言模型：从图文对齐到图文对话}

有了 CLIP / CoCa 这条线，下一步自然就是把视觉 token 接进语言模型。讲者先提到了 LLaVA，再提 Flamingo，说明多模态模型的真正目的，不是只做“看图分类”，而是让视觉成为 LLM 推理的上下文。\footnote{字幕 00:26:00--00:29:00 先把视觉语言模型和 autoregressive 语言模型联系起来，再引出 LLaVA / Flamingo。}

\subsection{LLaVA：最直接的连接方式}

LLaVA 的思路很朴素：先用 CLIP image encoder 把图像变成一组 token，再用一个小的投影层把这些视觉 token 映射到 LLM 能理解的空间，最后把它们和文本 token 拼在一起做 next-token prediction。它的关键不是“又发明了一套视觉编码”，而是把已有的视觉特征和已有的语言模型接到了一起。

\begin{figure}[H]
\centering
\includegraphics[width=0.95\textwidth]{slides-images/slide-032.jpg}
\caption{多模态语言模型的第一步：把视觉 token 接进语言模型}
\end{figure}

\begin{knowledgebox}{LLaVA 的关键组件}
\begin{itemize}
    \item 图像编码器负责把图像压成 token。
    \item 一个线性投影把视觉特征对齐到 LLM embedding 空间。
    \item LLM 负责生成回答、描述和推理。
\end{itemize}
\end{knowledgebox}

讲者还指出，CLIP encoder 并不只有 CLS token 真正有用，penultimate layer 的特征更有空间信息，因此在做多模态接入时，通常会偏向使用更底层、更细粒度的视觉表示，而不是只用最终聚合向量。\footnote{字幕 00:30:19--00:30:59 对 penultimate features 的用途有直接说明。}

\subsection{Flamingo：把视觉上下文喂进每一层}

Flamingo 是这讲里更重要的多模态架构案例。它延续了“视觉 encoder + LLM”的总体方向，但做法比简单拼接更精细。它把视觉特征先经过 perceiver sampler 压缩，再用 gated cross-attention 把视觉上下文注入 LLM 的每一层。也就是说，语言模型不是只在输入端看一眼图像，而是在每层都能选择性回看视觉信息。\footnote{字幕 00:31:20--00:35:10 对 Flamingo 的 cross-attention 和 perceiver sampler 有详细解释。}

\begin{figure}[H]
\centering
\includegraphics[width=0.95\textwidth]{slides-images/slide-034.jpg}
\caption{Flamingo 的核心思想：冻结大部分语言模型，只训练跨模态桥接模块}
\end{figure}

\begin{figure}[H]
\centering
\includegraphics[width=0.95\textwidth]{slides-images/slide-035.jpg}
\caption{Flamingo 通过 gated cross-attention 把视觉信息注入每一层语言模型}
\end{figure}

\begin{importantbox}{Flamingo 的设计重点}
\begin{itemize}
    \item 视觉 encoder 和语言模型大多冻结，训练成本更可控。
    \item cross-attention 让 LLM 在生成时能动态读取视觉上下文。
    \item perceiver sampler 把视觉 token 压到固定长度，避免输入过长。
\end{itemize}
\end{importantbox}

\subsection{多轮对话和 few-shot reasoning}

Flamingo 的强项不是单轮 caption，而是多轮对话与 in-context learning。你可以先给几个图文示例，再让它描述新的图；也可以先给几个问答示例，再问新图。它会把这些上下文当成“例子”，而不是重新训练一个新任务模型。\footnote{字幕 00:37:00--00:39:30 展示了多轮对话、few-shot captioning 和 question answering。}

\begin{figure}[H]
\centering
\includegraphics[width=0.95\textwidth]{slides-images/slide-037.jpg}
\caption{Flamingo 可以进行多轮图文对话，而不是只输出单句 caption}
\end{figure}

\begin{figure}[H]
\centering
\includegraphics[width=0.95\textwidth]{slides-images/slide-038.jpg}
\caption{Flamingo 的 few-shot / in-context 学习能力：给例子，就能模仿行为}
\end{figure}

\begin{knowledgebox}{为什么 in-context learning 重要}
它让多模态模型更像“通用接口”，而不是“专用分类器”。研究者不必为每个任务重新微调，只要通过示例就能让模型进入新的任务模式。
\end{knowledgebox}

\subsection{Grounding 到像素，才能减少幻觉}

讲者后来介绍了一个很有意思的点：如果模型回答“有几条船”，它很容易直接猜一个数字。但如果模型必须先“指向”图中被计数的对象，再给出最终答案，决策就更 grounded，幻觉也会少一些。Momo 就是这样一种做法：先把 decision-making grounded in the pixels，再做最终输出。\footnote{字幕 00:47:29--00:48:30 对 grounding 到 pixels 的设计做了直接说明。}

\begin{figure}[H]
\centering
\includegraphics[width=0.95\textwidth]{slides-images/slide-044.jpg}
\caption{让模型先指向像素，再做最终判断，是降低幻觉的一种方法}
\end{figure}

\begin{importantbox}{grounding 的意义}
如果模型必须先指出证据，再给出答案，它就更难完全靠语言先验胡编。grounding 不是万能解药，但它会迫使模型把输出和视觉证据绑定起来。
\end{importantbox}

\subsection{Momo：开放、可复现、可评估的多模态模型}

讲者还介绍了他们自己的 Momo。它的亮点不是“更大”，而是更开放：open weights、open data、open code，外加公开评测和用户研究。训练数据规模只有约 70 万图文对，但因为做了精心筛选和 grounding，效果能和更大模型接近。讲者甚至提到，他们做了 870 名用户、325,000 对 pairwise comparison 的用户研究，Momo 在 ELO 上逼近 GPT-4o。\footnote{字幕 00:47:32--00:49:00 和 00:50:00--00:50:40 讨论了 Momo 的数据规模、用户研究和 ELO 结果。}

\begin{figure}[H]
\centering
\includegraphics[width=0.95\textwidth]{slides-images/slide-048.jpg}
\caption{Momo 的一个关键特征：用更少但更精心筛选的数据，配合 grounding 学到强的多模态能力}
\end{figure}

\begin{knowledgebox}{为什么 open matters}
如果模型只靠私有数据和私有训练细节，社区就无法知道为什么它会强，也无法独立复现。开放 weights、data 和 code，才可能把模型能力真正扩散到研究社区。
\end{knowledgebox}

